\documentclass[11pt,a4paper]{article}
\usepackage[utf8]{inputenc}
\usepackage[T1]{fontenc}
\usepackage{geometry}
\geometry{margin=1.6cm}
\usepackage{booktabs}
\usepackage{longtable}
\usepackage{array}
\usepackage{graphicx}
\usepackage{tabularx}
\usepackage{fancyvrb}
\usepackage{xcolor}
\usepackage[hidelinks]{hyperref}

\newcommand{\pendiente}[1]{\textcolor{red}{\textbf{PENDIENTE:} #1}}

\title{Informe final del proyecto S-MIPS}
\author{}
\date{}

\begin{document}

\noindent \textbf{Nombre: Alejandro Zulueta Tamayo} \underline{\hspace{10cm}}\\[4mm]

% =====================================================================
\section{Identificación y componentes implementados en VHDL}

El circuito entregado \texttt{circuit/smips.circ} implementa la CPU multiciclo S-MIPS dentro del subcircuito \texttt{CPU}, usando exclusivamente componentes de Logisim Evolution (no se ha usado VHDL). Por tanto, \textbf{no hay archivos \texttt{.vhd} ni entidades VHDL incrustadas} en el proyecto.

\subsection{Estructura general implementada}

La placa (memoria, reloj, reset, terminal y LED de halt) se entrega resuelta. El trabajo realizado está dentro de \texttt{CPU}, que se compone de:

\begin{itemize}
  \item \texttt{Datapath}: contiene el banco de registros, los registros temporales (\texttt{PC}, \texttt{IR}, \texttt{A}, \texttt{B}, \texttt{BT}, \texttt{ALUOut}, \texttt{MDR}), la ALU, el extensor de inmediato, el manejador de bytes y los multiplexores de selección de datos.
  \item \texttt{Instruction\_Decoder}: decodifica los campos \texttt{OP} y \texttt{FUNCT} y genera las señales de clase de instrucción (\texttt{I\_*)} y las funciones (\texttt{F\_*}).
  \item \texttt{Branch\_Control}: calcula \texttt{BRANCH\_TAKEN} a partir de \texttt{I\_BNE}, \texttt{I\_BRANCH} y la señal \texttt{ZERO} de la ALU.
  \item \texttt{Control\_Unit}: FSM de 18 estados que genera todas las señales de control.
\end{itemize}

% =====================================================================
\section{Diagrama de la ruta de datos con señales de control}

La Figura~\ref{fig:datapath} muestra la ruta de datos simplificada. Las señales entre corchetes son las señales de control generadas por \texttt{Control\_Unit}.

\begin{figure}[htbp]
\centering
\begin{Verbatim}[fontsize=\small]
                    [PC_WRITE, PC_SRC]
                            |
  +------+    +-------------v------+    +--------+    +-------------------+
  |  PC  |--->| Memoria de Instr.  |--->|   IR   |--->| Instruction_Decoder|
  +------+    +--------------------+    +--------+    +---------+---------+
     ^                                     |                    |
     |                                     |                    v
     |                              [IR_WRITE]         I_*, F_*, REG_DST,
     |                                                 MEM_TO_REG, EXT_OP
     |                                                      |
     |                    +---------+    +----------+        v
     +--------------------|  ALU    |<---|   A/B    |<--- Banco de registros
     |   [PC_SRC=0]       +----+----+    +----+-----+     [A_WRITE, B_WRITE]
     |                        |             ^
     |                  [ALU_CONTROL,        |
     |                   ALU_SRC_A/B]        |
     |                        v             |
     |                  +-----------+        |
     +------------------|  ALUOut   |--------+
                        +-----+-----+
                              | [ALU_OUT_WRITE]
                              v
                        +-----------+     [IOR_D]
                        |    MDR    |-------------> Memoria de Datos
                        +-----+-----+              [MEM_WE, MEM_BE]
                              | [MDR_WRITE]
                              v
                        +-----------+     [MEM_BYTE, I_LB]
                        |Byte_Handler|------------> MEM_TO_REG -> REG_WRITE
                        +-----------+
\end{Verbatim}
\caption{Ruta de datos simplificada de la CPU S-MIPS. Entre corchetes, señales de control.}
\label{fig:datapath}
\end{figure}

\paragraph{Señales de control de la ruta de datos.}
Las señales generadas por \texttt{Control\_Unit} son:

\begin{itemize}
  \item \textbf{Escrituras:} \texttt{REG\_WRITE}, \texttt{IR\_WRITE}, \texttt{PC\_WRITE}, \texttt{A\_WRITE}, \texttt{B\_WRITE}, \texttt{BT\_WRITE}, \texttt{MDR\_WRITE}, \texttt{ALU\_OUT\_WRITE}, \texttt{MEM\_WE}, \texttt{TTY\_WE}, \texttt{HALT}.
  \item \textbf{Selección:} \texttt{PC\_SRC} (2 bits), \texttt{ALU\_SRC\_A} (2 bits), \texttt{ALU\_SRC\_B} (2 bits), \texttt{IOR\_D}, \texttt{ALU\_CONTROL} (3 bits), \texttt{EXT\_OP} (2 bits), \texttt{REG\_DST} (2 bits), \texttt{MEM\_TO\_REG} (2 bits).
\end{itemize}

El multiplexor de dirección de memoria usa \texttt{IOR\_D}: con \texttt{IOR\_D=0} se accede a instrucciones (dirección =\ \texttt{PC}); con \texttt{IOR\_D=1} se accede a datos (dirección =\ \texttt{ALUOut}, calculada como base+offset en \texttt{S\_MEM\_ADDR}).

% =====================================================================
\section{Máquina de estados: estados, transiciones y condiciones}

La unidad de control implementa una FSM de 18 estados. Cada estado tiene un registro de 1 bit: \texttt{S\_*} es la salida \texttt{Q} (estado actual) y \texttt{N\_*} es la entrada \texttt{D} (próximo estado). En cada flanco de subida de \texttt{CLK}, \texttt{S\_X(t+1) <- N\_X(t)}. Todos los registros están siempre habilitados con \texttt{ONE1} y se reinician con \texttt{RESET}.

\subsection{Transiciones}

La Tabla~\ref{tab:fsm} recoge todos los estados y sus transiciones.

\begin{longtable}{@{}p{2.7cm}p{4.2cm}p{8.5cm}@{}}
\caption{Estados y transiciones de la FSM.}
\label{tab:fsm}\\
\toprule
Estado actual & Condición & Próximo estado \\
\midrule
\endfirsthead
\toprule
Estado actual & Condición & Próximo estado \\
\midrule
\endhead
\texttt{S\_FETCH} & siempre & \texttt{S\_DECODE} \\
\texttt{S\_DECODE} & \texttt{I\_ALU\_R} & \texttt{S\_EXEC\_R} \\
\texttt{S\_DECODE} & \texttt{I\_I\_ALU} & \texttt{S\_EXEC\_I} \\
\texttt{S\_DECODE} & \texttt{I\_LOAD} $\lor$ \texttt{I\_STORE} & \texttt{S\_MEM\_ADDR} \\
\texttt{S\_DECODE} & \texttt{I\_BRANCH} & \texttt{S\_BRANCH} \\
\texttt{S\_DECODE} & \texttt{I\_J} & \texttt{S\_JUMP} \\
\texttt{S\_DECODE} & \texttt{I\_JAL} & \texttt{S\_JAL} \\
\texttt{S\_DECODE} & \texttt{I\_JR} & \texttt{S\_JR} \\
\texttt{S\_DECODE} & \texttt{I\_SYSCALL} & \texttt{S\_SYSCALL} \\
\texttt{S\_DECODE} & \texttt{ILLEGAL} & \texttt{S\_HALT} \\
\texttt{S\_EXEC\_R} & siempre & \texttt{S\_WB\_R} \\
\texttt{S\_WB\_R} & siempre & \texttt{S\_FETCH} \\
\texttt{S\_EXEC\_I} & siempre & \texttt{S\_WB\_I} \\
\texttt{S\_WB\_I} & siempre & \texttt{S\_FETCH} \\
\texttt{S\_MEM\_ADDR} & \texttt{I\_LOAD} & \texttt{S\_MEM\_RD} \\
\texttt{S\_MEM\_ADDR} & \texttt{I\_STORE} & \texttt{S\_MEM\_WR} \\
\texttt{S\_MEM\_RD} & siempre & \texttt{S\_MEM\_WB} \\
\texttt{S\_MEM\_WB} & siempre & \texttt{S\_FETCH} \\
\texttt{S\_MEM\_WR} & siempre & \texttt{S\_FETCH} \\
\texttt{S\_BRANCH} & \texttt{BRANCH\_TAKEN}=1 & \texttt{S\_BRANCH\_T} \\
\texttt{S\_BRANCH} & \texttt{BRANCH\_TAKEN}=0 & \texttt{S\_FETCH} \\
\texttt{S\_BRANCH\_T} & siempre & \texttt{S\_FETCH} \\
\texttt{S\_JUMP} & siempre & \texttt{S\_FETCH} \\
\texttt{S\_JAL} & siempre & \texttt{S\_FETCH} \\
\texttt{S\_JR} & siempre & \texttt{S\_FETCH} \\
\texttt{S\_SYSCALL} & \texttt{V0}=10 & \texttt{S\_HALT} \\
\texttt{S\_SYSCALL} & \texttt{V0}=11 & \texttt{S\_TTY} \\
\texttt{S\_SYSCALL} & otro \texttt{V0} & \texttt{S\_FETCH} \\
\texttt{S\_TTY} & siempre & \texttt{S\_FETCH} \\
\texttt{S\_HALT} & siempre & \texttt{S\_HALT} \\
\bottomrule
\end{longtable}

\subsection{Detalle de la lógica de próximo estado}

\begin{itemize}
  \item \textbf{Arranque:} tras \texttt{RESET}, todos los \texttt{S\_*} valen 0; \texttt{ANY\_STATE}=0 y \texttt{NO\_ACTIVE\_STATE}=1, por lo que \texttt{N\_S\_FETCH} se activa y la FSM arranca en \texttt{S\_FETCH}.
  \item \textbf{Retorno a fetch:} \texttt{N\_S\_FETCH} es el OR de \texttt{S\_WB\_R}, \texttt{S\_WB\_I}, \texttt{S\_MEM\_WB}, \texttt{S\_MEM\_WR}, \texttt{S\_BRANCH\_T}, \texttt{T\_BRANCH\_NT} (=\ \texttt{S\_BRANCH} $\land$ $\neg$\texttt{BRANCH\_TAKEN}), \texttt{S\_JUMP}, \texttt{S\_JAL}, \texttt{S\_JR}, \texttt{S\_TTY}, \texttt{T\_SYS\_CONT} (syscall con \texttt{V0}$\neq$10 y \texttt{V0}$\neq$11) y \texttt{NO\_ACTIVE\_STATE}.
  \item \textbf{Syscall:} el valor de \texttt{V0} se compara con las constantes \texttt{SC\_EXIT\_10}=10 y \texttt{SC\_PUTC\_11}=11. \texttt{V0\_EQ\_11} entrega \texttt{N\_S\_TTY} = \texttt{S\_SYSCALL} $\land$ \texttt{V0\_EQ\_11}; \texttt{V0\_EQ\_10} contribuye a \texttt{N\_S\_HALT}.
  \item \textbf{Halt:} \texttt{N\_S\_HALT} = \texttt{S\_HALT} $\lor$ \texttt{N\_HALT\_ILL} $\lor$ \texttt{N\_HALT\_SYS}, por lo que una vez en halt el estado es absorbente.
\end{itemize}

% =====================================================================
\section{Tabla de señales de control por estado}

La Tabla~\ref{tab:controls} muestra el valor de cada señal de control en cada estado. `1' indica activo, `0' inactivo, `X' indica ``no importa'' (la señal no es relevante en ese estado), `D' indica que se toma del decodificador y `C' que se fuerza a una constante.

\begin{longtable}{@{}p{2.25cm}cccccccccccc@{}}
\caption{Valores de las señales de control en cada estado.}
\label{tab:controls}\\
\toprule
Estado & \rotatebox{90}{REG\_WRITE} & \rotatebox{90}{IR\_WRITE} & \rotatebox{90}{PC\_WRITE} & \rotatebox{90}{A/B\_WRITE} & \rotatebox{90}{ALU\_OUT\_W} & \rotatebox{90}{MDR\_WRITE} & \rotatebox{90}{MEM\_WE} & \rotatebox{90}{TTY\_WE} & \rotatebox{90}{IOR\_D} & \rotatebox{90}{PC\_SRC} & \rotatebox{90}{ALU\_SRC\_A/B} & \rotatebox{90}{ALU\_CTRL} \\
\midrule
\endfirsthead
\toprule
Estado & \rotatebox{90}{REG\_WRITE} & \rotatebox{90}{IR\_WRITE} & \rotatebox{90}{PC\_WRITE} & \rotatebox{90}{A/B\_WRITE} & \rotatebox{90}{ALU\_OUT\_W} & \rotatebox{90}{MDR\_WRITE} & \rotatebox{90}{MEM\_WE} & \rotatebox{90}{TTY\_WE} & \rotatebox{90}{IOR\_D} & \rotatebox{90}{PC\_SRC} & \rotatebox{90}{ALU\_SRC\_A/B} & \rotatebox{90}{ALU\_CTRL} \\
\midrule
\endhead
\texttt{S\_FETCH}    & 0 & 1 & 1 & 0 & 0 & 0 & 0 & 0 & 0 & 0 & PC/4 & ADD \\
\texttt{S\_DECODE}  & 0 & 0 & 0 & 1 & 0 & 0 & 0 & 0 & 0 & X & PC/X & ADD \\
\texttt{S\_EXEC\_R} & 0 & 0 & 0 & 0 & 1 & 0 & 0 & 0 & 0 & X & A/B $\lor$ shamt & D \\
\texttt{S\_WB\_R}   & 1 & 0 & 0 & 0 & 0 & 0 & 0 & 0 & 0 & X & X & X \\
\texttt{S\_EXEC\_I} & 0 & 0 & 0 & 0 & 1 & 0 & 0 & 0 & 0 & X & A/Imm & D \\
\texttt{S\_WB\_I}   & 1 & 0 & 0 & 0 & 0 & 0 & 0 & 0 & 0 & X & X & X \\
\texttt{S\_MEM\_ADDR} & 0 & 0 & 0 & 0 & 1 & 0 & 0 & 0 & 0 & X & A/Imm & ADD \\
\texttt{S\_MEM\_RD}  & 0 & 0 & 0 & 0 & 0 & 1 & 0 & 0 & 1 & X & X & X \\
\texttt{S\_MEM\_WB}  & 1 & 0 & 0 & 0 & 0 & 0 & 0 & 0 & 1 & X & X & X \\
\texttt{S\_MEM\_WR}  & 0 & 0 & 0 & 0 & 0 & 0 & 1 & 0 & 1 & X & A/Imm & ADD \\
\texttt{S\_BRANCH}   & 0 & 0 & 0 & 0 & 0 & 0 & 0 & 0 & 0 & X & A/B & SUB \\
\texttt{S\_BRANCH\_T}& 0 & 0 & 1 & 0 & 0 & 0 & 0 & 0 & 0 & 1 & X & X \\
\texttt{S\_JUMP}     & 0 & 0 & 1 & 0 & 0 & 0 & 0 & 0 & 0 & 2 & X & X \\
\texttt{S\_JAL}      & 1 & 0 & 1 & 0 & 0 & 0 & 0 & 0 & 0 & 2 & X & X \\
\texttt{S\_JR}       & 0 & 0 & 1 & 0 & 0 & 0 & 0 & 0 & 0 & 3 & X & X \\
\texttt{S\_SYSCALL}  & 0 & 0 & 0 & 0 & 0 & 0 & 0 & 0 & 0 & X & X & X \\
\texttt{S\_TTY}      & 0 & 0 & 0 & 0 & 0 & 0 & 0 & 1 & 0 & X & X & X \\
\texttt{S\_HALT}     & 0 & 0 & 0 & 0 & 0 & 0 & 0 & 0 & 0 & X & X & X \\
\bottomrule
\end{longtable}

\paragraph{Señales derivadas del estado.}
Además de la tabla, el control genera estos términos internos:

\begin{itemize}
  \item \texttt{T\_LS} = \texttt{I\_LOAD} $\lor$ \texttt{I\_STORE}.
  \item \texttt{T\_USE\_A} = \texttt{S\_EXEC\_I} $\lor$ \texttt{S\_MEM\_ADDR} $\lor$ \texttt{S\_BRANCH}.
  \item \texttt{T\_USE\_PC} = \texttt{S\_FETCH} $\lor$ \texttt{S\_DECODE}.
  \item \texttt{T\_USE\_IMM} = \texttt{S\_EXEC\_I} $\lor$ \texttt{S\_MEM\_ADDR}.
  \item \texttt{T\_ALU\_ADD} = \texttt{S\_FETCH} $\lor$ \texttt{S\_DECODE} $\lor$ \texttt{S\_MEM\_ADDR}.
  \item \texttt{T\_EXT\_SIGN} = \texttt{S\_DECODE} $\lor$ \texttt{S\_MEM\_ADDR} $\lor$ \texttt{S\_BRANCH}.
  \item \texttt{SEL\_PC\_JUMP\_JAL} = \texttt{S\_JUMP} $\lor$ \texttt{S\_JAL}.
\end{itemize}

\paragraph{Cascadas de multiplexores del control.}
Las salidas de selección se construyen con multiplexores 2:1 encadenados, donde la salida de una etapa entra por \texttt{d0} de la siguiente:

\begin{itemize}
  \item \texttt{PC\_SRC}: \texttt{PC\_01 = S\_BRANCH\_T?1:0} $\rightarrow$ \texttt{PC\_012 = SEL\_PC\_JUMP\_JAL?2:PC\_01} $\rightarrow$ \texttt{PC\_SRC = S\_JR?3:PC\_012}.
  \item \texttt{ALU\_SRC\_A}: \texttt{ASA\_01 = T\_SEL\_A?1:0} $\rightarrow$ \texttt{ASA\_012 = T\_SHIFT?2:ASA\_01} $\rightarrow$ \texttt{ALU\_SRC\_A = T\_USE\_PC?0:ASA\_012}.
  \item \texttt{ALU\_SRC\_B}: \texttt{ASB\_01 = S\_FETCH?1:0} $\rightarrow$ \texttt{ASB\_012 = T\_USE\_IMM?2:ASB\_01} $\rightarrow$ \texttt{ALU\_SRC\_B = S\_DECODE?3:ASB\_012}.
  \item \texttt{ALU\_CONTROL}: \texttt{AC\_1 = T\_ALU\_ADD?000:DEC\_ALU\_CONTROL} $\rightarrow$ \texttt{ALU\_CONTROL = S\_BRANCH?001:AC\_1}.
\end{itemize}

\paragraph{Selector directo:}
\texttt{EXT\_OP = T\_EXT\_SIGN?0:DEC\_EXT\_OP}; \texttt{REG\_DST = S\_JAL?2:DEC\_REG\_DST}; \texttt{MEM\_TO\_REG = S\_JAL?2:DEC\_MEM\_TO\_REG}.

% =====================================================================
\section{Codificación de operaciones de la ALU y salidas del decodificador}

\subsection{Codificación de la ALU}

La ALU usa \texttt{ALU\_CONTROL} de 3 bits con un decoder 3:8 y un multiplexor 8:1. La Tabla~\ref{tab:alu} recoge la codificación implementada.

\begin{table}[htbp]
\centering
\caption{Codificación de \texttt{ALU\_CONTROL}.}
\label{tab:alu}
\begin{tabular}{@{}cll@{}}
\toprule
Código & Operación & Uso típico \\
\midrule
000 & \texttt{A+B} (ADD) & PC+4, cálculo de direcciones, add/addu/addi/addiu \\
001 & \texttt{A$-$B} (SUB) & sub, comparación de ramas (bit~0 invierte B y activa \texttt{Cin}) \\
010 & \texttt{A $\land$ B} & and/andi \\
011 & \texttt{A $\lor$ B} & or/ori \\
100 & \texttt{A $<$ B ? 1 : 0} (SLT) & slt/slti \\
101 & \texttt{A $\ll$ shamt} & sll \\
110 & \texttt{A $\gg$ shamt} lógico & srl \\
111 & \texttt{B $\ll$ 16} (LUI) & lui \\
\bottomrule
\end{tabular}
\end{table}

\noindent La codificación se obtiene del \texttt{Decoder\_3} y del \texttt{Multiplexer\_3} de \texttt{ALU}: el código~0 y el~1 seleccionan la misma salida del sumador, pero el bit~0 además invierte $B$ y activa el acarreo de entrada, con lo que el código~1 realiza la resta. El bit \texttt{ZERO} sale de un comparador de 32 bits entre el resultado y cero; lo usa \texttt{Branch\_Control} junto con \texttt{I\_BNE} para calcular \texttt{BRANCH\_TAKEN}.

\subsection{Salidas del decodificador de instrucción}

El \texttt{Instruction\_Decoder} recibe \texttt{OP} (6 bits) y \texttt{FUNCT} (6 bits). Genera señales de clase \texttt{I\_*} y de función \texttt{F\_*}. Los valores de \texttt{OP} y \texttt{FUNCT} se escriben en hexadecimal de dos dígitos: \texttt{OP=00} equivale a \texttt{0x00}, \texttt{FUNCT=2A} equivale a \texttt{0x2A}, etc. La Tabla~\ref{tab:decoder} resume la decodificación.

\begin{longtable}{@{}p{2.4cm}p{2.2cm}p{4.5cm}p{6.0cm}@{}}
\caption{Salidas del decodificador.}
\label{tab:decoder}\\
\toprule
Señal & Tipo & Condición & Uso \\
\midrule
\endfirsthead
\toprule
Señal & Tipo & Condición & Uso \\
\midrule
\endhead
\texttt{I\_ALU\_R} & clase & \texttt{OP=00} y \texttt{FUNCT} $\in$ \{00, 02, 20, 21, 22, 24, 25, 2A\} & ejecuta R-type \\
\texttt{I\_I\_ALU} & clase & \texttt{OP=08(addi), 09(addiu), 0A(slti), 0C(andi), 0D(ori)} & ejecuta I-type ALU \\
\texttt{I\_LOAD} & clase & \texttt{OP=20(lb), 23(lw), 24(lbu)} & acceso a memoria de lectura \\
\texttt{I\_STORE} & clase & \texttt{OP=28(sb), 2B(sw)} & acceso a memoria de escritura \\
\texttt{I\_BRANCH} & clase & \texttt{OP=04(beq), 05(bne)} & evaluación de rama \\
\texttt{I\_J} & instr. & \texttt{OP=02} & salto \texttt{j} \\
\texttt{I\_JAL} & instr. & \texttt{OP=03} & salto y link \\
\texttt{I\_JR} & instr. & \texttt{OP=00} $\land$ \texttt{FUNCT=08} & salto registrado \\
\texttt{I\_SYSCALL} & instr. & \texttt{OP=00} $\land$ \texttt{FUNCT=0C} & llamada al sistema \\
\texttt{I\_SHIFT} & clase & \texttt{OP=00} $\land$ \texttt{FUNCT=00(sll), 02(srl)} & selecciona \texttt{shamt} en ALU \\
\texttt{ILLEGAL} & estado & ninguna combinación válida de \texttt{OP}/\texttt{FUNCT} & halt \\
\texttt{MEM\_BYTE} & dato & \texttt{OP=20(lb), 24(lbu), 28(sb)} & selecciona byte en \texttt{Byte\_Handler} \\
\texttt{F\_ADD} & FUNCT & \texttt{OP=00} $\land$ \texttt{FUNCT=20} (add) & operación ALU R-type \\
\texttt{F\_ADDU} & FUNCT & \texttt{OP=00} $\land$ \texttt{FUNCT=21} (addu) & operación ALU R-type \\
\texttt{F\_SUB} & FUNCT & \texttt{OP=00} $\land$ \texttt{FUNCT=22} (sub) & operación ALU R-type \\
\texttt{F\_AND} & FUNCT & \texttt{OP=00} $\land$ \texttt{FUNCT=24} (and) & operación ALU R-type \\
\texttt{F\_OR} & FUNCT & \texttt{OP=00} $\land$ \texttt{FUNCT=25} (or) & operación ALU R-type \\
\texttt{F\_SLT} & FUNCT & \texttt{OP=00} $\land$ \texttt{FUNCT=2A} (slt) & operación ALU R-type \\
\texttt{F\_SLL} & FUNCT & \texttt{OP=00} $\land$ \texttt{FUNCT=00} (sll) & desplazamiento \\
\texttt{F\_SRL} & FUNCT & \texttt{OP=00} $\land$ \texttt{FUNCT=02} (srl) & desplazamiento \\
\texttt{F\_JR} & FUNCT & \texttt{OP=00} $\land$ \texttt{FUNCT=08} (jr) & salto registrado \\
\texttt{F\_SYSCALL} & FUNCT & \texttt{OP=00} $\land$ \texttt{FUNCT=0C} (syscall) & llamada al sistema \\
\bottomrule
\end{longtable}

Las señales \texttt{I\_SUB}, \texttt{I\_AND}, \texttt{I\_OR}, \texttt{I\_SLT}, \texttt{I\_SLL} e \texttt{I\_SRL} usadas en el apartado~\ref{sec:codigos} son las funciones de la tabla~\ref{tab:decoder} después de aplicar \texttt{OP=00}; por ejemplo, \texttt{I\_SUB = OP=00} $\land$ \texttt{F\_SUB} y \texttt{I\_AND = OP=00} $\land$ \texttt{F\_AND}. Por eso en la tabla aparecen como \texttt{F\_*} con \texttt{OP=00} explícito en la condición.

\subsection{Códigos propuestos por el decodificador}
\label{sec:codigos}

Los buses \texttt{ALU\_CONTROL}, \texttt{EXT\_OP}, \texttt{REG\_DST} y \texttt{MEM\_TO\_REG} se forman en el propio \texttt{Instruction\_Decoder} combinando señales \texttt{I\_*}/\texttt{F\_*} con compuertas OR y un \emph{splitter} usado como combinador de bits. La Control\_Unit puede después sustituirlos por constantes según el estado (tabla~\ref{tab:controls}).

\paragraph{\texttt{ALU\_CONTROL} (3 bits).} Cada bit es el OR de las instrucciones que requieren ese bit a~1:

\begin{itemize}
  \item bit~0 = \texttt{I\_SUB} $\lor$ \texttt{I\_BEQ} $\lor$ \texttt{I\_BNE} $\lor$ \texttt{I\_OR} $\lor$ \texttt{I\_ORI} $\lor$ \texttt{I\_SLL} $\lor$ \texttt{I\_LUI}
  \item bit~1 = \texttt{I\_AND} $\lor$ \texttt{I\_ANDI} $\lor$ \texttt{I\_OR} $\lor$ \texttt{I\_ORI} $\lor$ \texttt{I\_SRL} $\lor$ \texttt{I\_LUI}
  \item bit~2 = \texttt{I\_SLT} $\lor$ \texttt{I\_SLTI} $\lor$ \texttt{I\_SLL} $\lor$ \texttt{I\_SRL} $\lor$ \texttt{I\_LUI}
\end{itemize}

Por ejemplo, \texttt{add} produce 000 (ADD), \texttt{sub} 001 (SUB), \texttt{and} 010, \texttt{or} 011, \texttt{slt} 100, \texttt{sll} 101, \texttt{srl} 110 y \texttt{lui} 111, coherente con la tabla~\ref{tab:alu}.

\paragraph{\texttt{EXT\_OP} (2 bits).} bit~0 = \texttt{I\_ANDI} $\lor$ \texttt{I\_ORI}; bit~1 = \texttt{I\_LUI}. El bit~0 selecciona la extensión con ceros para inmediatos lógicos, el bit~1 la colocación del inmediato en los 16 bits altos para \texttt{lui}; con ambos a~0 se usa extensión con signo.

\paragraph{\texttt{REG\_DST} (2 bits).} bit~0 = \texttt{I\_ALU\_R}; bit~1 = \texttt{I\_JAL}. Selecciona \texttt{IR[rt]}, \texttt{IR[rd]} o el registro~31 ($ra$) según la clase de instrucción.

\paragraph{\texttt{MEM\_TO\_REG} (2 bits).} bit~0 = \texttt{I\_LOAD}; bit~1 = \texttt{I\_JAL}. Selecciona el resultado de la ALU, el dato leído (\texttt{MDR}) o el \texttt{PC} de retorno para escribir en el banco de registros.

% =====================================================================
\section{Número de ciclos de cada clase de instrucción}

Contando desde que se entra en \texttt{S\_FETCH} hasta que se vuelve a él (incluido el ciclo de fetch):

\begin{table}[htbp]
\centering
\caption{Ciclos por clase de instrucción.}
\label{tab:ciclos}
\begin{tabular}{@{}llc@{}}
\toprule
Clase & Recorrido & Ciclos \\
\midrule
R-type ALU & FETCH $\to$ DECODE $\to$ EXEC\_R $\to$ WB\_R & 4 \\
I-type ALU & FETCH $\to$ DECODE $\to$ EXEC\_I $\to$ WB\_I & 4 \\
Load (\texttt{lw/lb/lbu}) & FETCH $\to$ DECODE $\to$ MEM\_ADDR $\to$ MEM\_RD $\to$ MEM\_WB & 5 \\
Store (\texttt{sw/sb}) & FETCH $\to$ DECODE $\to$ MEM\_ADDR $\to$ MEM\_WR & 4 \\
Rama (\texttt{beq/bne}) & FETCH $\to$ DECODE $\to$ BRANCH [$\to$ BRANCH\_T] & 3 o 4 \\
Salto (\texttt{j/jal/jr}) & FETCH $\to$ DECODE $\to$ JUMP/JAL/JR & 3 \\
Syscall 11 & FETCH $\to$ DECODE $\to$ SYSCALL $\to$ TTY & 4 \\
Syscall 10 & FETCH $\to$ DECODE $\to$ SYSCALL $\to$ HALT & 4 \\
\bottomrule
\end{tabular}
\end{table}

% =====================================================================
\section{Cobertura de los programas de \texttt{tests/}}

La Tabla~\ref{tab:cobertura} indica en qué programa se cubre cada fila de la tabla de casos de la guía (R21).

\begin{longtable}{@{}p{4.6cm}p{4.0cm}p{6.5cm}@{}}
\caption{Cobertura de los casos obligatorios por los programas de \texttt{tests/}.}
\label{tab:cobertura}\\
\toprule
Filas de la tabla de casos & Programa(s) & Casos cubiertos \\
\midrule
\endfirsthead
\toprule
Filas de la tabla de casos & Programa(s) & Casos cubiertos \\
\midrule
\endhead
\texttt{add}, \texttt{sub} & \texttt{01\_arithmetic.asm} & positivos y con operando negativo \\
\texttt{and}, \texttt{or} & \texttt{02\_logic.asm} & operandos con unos en los 16 bits altos \\
\texttt{slt}, \texttt{slti} & \texttt{03\_comparisons.asm} & resultados 0 y 1; comparación con/sin signo con operando negativo \\
\texttt{sll}, \texttt{srl} & \texttt{04\_shifts.asm} & \texttt{shamt}=0 y $\neq$0; \texttt{srl} sobre negativo \\
\texttt{addi} & \texttt{01\_arithmetic.asm} & inmediato negativo \\
\texttt{addiu}, \texttt{addu} & \texttt{01\_arithmetic.asm} & vía \texttt{li} y \texttt{move} \\
\texttt{andi}, \texttt{ori} & \texttt{02\_logic.asm} & inmediato con bit 15=1, extensión con ceros \\
\texttt{lui} & \texttt{02\_logic.asm} & inmediato con bit 15=1 \\
\texttt{lw}, \texttt{sw} & \texttt{05\_word\_memory.asm} & desplazamiento positivo y negativo \\
\texttt{lb}, \texttt{lbu} & \texttt{06\_byte\_memory.asm} & byte con bit 7=1 en posiciones $d[1{:}0]$ distintas \\
\texttt{sb} & \texttt{06\_byte\_memory.asm} & cuatro posiciones, comprobado con \texttt{lw} \\
\texttt{beq}, \texttt{bne} & \texttt{07\_control\_flow.asm} & tomado y no tomado; salto hacia atrás \\
\texttt{j} & \texttt{07\_control\_flow.asm} & al menos una vez \\
\texttt{jal}, \texttt{jr} & \texttt{07\_control\_flow.asm} & procedimiento no hoja con pila y retorno \\
\texttt{syscall} & \texttt{07\_control\_flow.asm} & servicios 11 y 10 \\
\texttt{\$zero} & \texttt{07\_control\_flow.asm} & escritura en \$zero y comprobación posterior \\
integración & \texttt{08\_integration.asm} & recurrencia, lógica, comparaciones, desplazamientos, memoria de palabra y byte, bucles anidados, pila y \texttt{syscall} \\
\bottomrule
\end{longtable}

% =====================================================================
\section{Resultados de los tests}

El archivo \texttt{salida-tests.txt} contiene la salida completa de \texttt{python3 tools/run\_tests.py} sobre la versión entregada. El resultado es \textbf{15/15 programas correctos}, incluyendo los 7 programas de \texttt{programs/} y los 8 de \texttt{tests/}:

\begin{Verbatim}[fontsize=\small]
Ejecutando 15 programas en la placa...
OK    programs/arithmetic.asm  (6664 ciclos)
OK    programs/array_sum.asm  (915 ciclos)
OK    programs/branches.asm  (926 ciclos)
OK    programs/gcd.asm  (2851 ciclos)
OK    programs/hello.asm  (297 ciclos)
OK    programs/memory.asm  (5075 ciclos)
OK    programs/multiply.asm  (1971 ciclos)
OK    tests/01_arithmetic.asm  (2103 ciclos)
OK    tests/02_logic.asm  (3009 ciclos)
OK    tests/03_comparisons.asm  (3061 ciclos)
OK    tests/04_shifts.asm  (3071 ciclos)
OK    tests/05_word_memory.asm  (2374 ciclos)
OK    tests/06_byte_memory.asm  (3696 ciclos)
OK    tests/07_control_flow.asm  (2933 ciclos)
OK    tests/08_integration.asm  (5059 ciclos)

15/15 correctos
\end{Verbatim}

\end{document}
